\section{Starobinsky Inflation from the $\sigma$-Field Action}
\label{sec:inflation_from_sigma}

The QGD action contains the $\sigma$-kinetic term
\begin{equation}
S_\sigma = \frac{\hbar^2}{2M}\int d^4x\sqrt{-g}\, g^{\mu\nu}\nabla_\mu\sigma^\alpha\nabla_\nu\sigma_\alpha,
\end{equation}
where the metric is composite: $g_{\mu\nu}=\eta_{\mu\nu}-\sigma_\mu\sigma_\nu$ (in the local Lorentz frame) and then transformed to FRW coordinates via $T^\alpha{}_\mu=\operatorname{diag}(1,a,a,a)$. For a homogeneous and isotropic universe we set $\sigma_\mu=(\varphi(t),0,0,0)$. The metric then becomes
\begin{equation}
ds^2 = (1-\varphi^2)dt^2 - a(t)^2 d\mathbf{x}^2.
\end{equation}
To obtain the standard FRW form with proper time we introduce a new time coordinate $\tau$ by $d\tau = \sqrt{1-\varphi^2}\,dt$. The metric becomes $ds^2 = -d\tau^2 + a^2 d\mathbf{x}^2$ and the scale factor $a$ is now a function of $\tau$. The $\sigma$-field is related to the lapse $N=1/\sqrt{1-\varphi^2}$. The kinetic term turns into
\begin{equation}
S_\sigma = \frac{\hbar^2}{2M}\int d\tau\, a^3\, \left(\frac{d\varphi}{d\tau}\right)^{\!2}\frac{1}{1-\varphi^2}.
\end{equation}
In terms of the lapse, $\varphi=\sqrt{1-1/N^2}$, so the kinetic term becomes a function of $N$ and $\dot N$.

\subsection{Fourth‑order dynamics and the massive mode}
The full $\sigma$-field equation derived in Chapter~5 includes a fourth‑order quantum stiffness term:
\begin{equation}
\ddot\varphi + 3H\dot\varphi - \kappa\ell_Q^2\!\left(\ddddot\varphi + 3H\dddot\varphi + \cdots\right) = 2\varphi\dot\varphi^2,
\label{eq:sigma_fourth}
\end{equation}
where $\ell_Q^2 = G\hbar/c^3$ and $\kappa\sim2$. In a slow‑roll inflationary regime the higher derivatives are small, but the stiffness term introduces a massive mode with mass $m_Q = \Mpl/\sqrt{\kappa}$ (the Planck mass). To see this, perturb around a de Sitter background $\varphi=\varphi_0+\delta\varphi$. In Fourier space the linearised equation gives the dispersion relation
\begin{equation}
-\omega^2\delta\varphi -3iH\omega\delta\varphi + \kappa\ell_Q^2\,\omega^4\delta\varphi = 0,
\end{equation}
which for high frequencies ($\omega\gg H$) yields $\omega^2\approx m_Q^2$. Thus the theory contains a Planck‑scale massive mode that can act as the inflaton.

\subsection{Mapping to the Starobinsky model}
Starobinsky inflation is described by the action
\begin{equation}
S = \int d^4x\sqrt{-g}\left(\frac{R}{16\pi G} + \alpha R^2\right).
\end{equation}
After a conformal transformation this becomes a scalar‑tensor theory with a scalar field $\phi$ (the scalaron) and potential
\begin{equation}
V(\phi) = \frac{M_{\mathrm{Pl}}^2}{8\alpha}\left(1 - e^{-\sqrt{2/3}\,\phi/M_{\mathrm{Pl}}}\right)^2.
\end{equation}
The spectral index and tensor‑to‑scalar ratio are
\begin{equation}
n_s = 1 - \frac{2}{N},\qquad r = \frac{12}{N^2},
\end{equation}
where $N$ is the number of e‑folds before the end of inflation. For $N=60$, $n_s\approx0.967$ and $r\approx0.003$, in excellent agreement with Planck~2018 data.

The fourth‑order term in QGD generates an effective $R^2$ term. The coefficient $\alpha$ can be estimated by matching the mass of the massive mode to the scalaron mass $m_\phi = 1/\sqrt{6\alpha}$. With $m_Q = \Mpl/\sqrt{\kappa}$ we obtain
\begin{equation}
\alpha = \frac{\kappa}{6}\,\ell_Q^2,
\end{equation}
since $\ell_Q = 1/\Mpl$ in natural units. In terms of the gravitational coupling,
\begin{equation}
\alpha \sim \frac{\kappa\ell_Q^2}{16\pi G} = \frac{\kappa}{16\pi G\Mpl^2} = \frac{\kappa}{16\pi}\,\ell_Q^2,
\end{equation}
which is indeed of order $\ell_Q^2$ up to numerical factors – the expected scale for an $R^2$ term.

\subsection{Verification of the spectral indices}
To confirm that the QGD dynamics reproduce the Starobinsky predictions we analyse the slow‑roll regime. In this regime the field equation~\eqref{eq:sigma_fourth} reduces to
\begin{equation}
3H\dot\varphi \approx 2\varphi\dot\varphi^2.
\end{equation}
From the Friedmann equation,
\begin{equation}
H^2 = \frac{4\pi G\hbar^2}{3c^2 M}\,\dot\varphi^2.
\end{equation}
Eliminating $\dot\varphi$ gives a relation between $H$ and $\varphi$. The slow‑roll parameters are
\begin{equation}
\epsilon = -\frac{\dot H}{H^2},\qquad \eta = \frac{\dot\epsilon}{\epsilon H}.
\end{equation}
Using the field equation one can express these in terms of $\varphi$. The number of e‑folds is
\begin{align}
N &= \int H\,dt = \int\frac{H}{\dot\varphi}\,d\varphi
   = \int\frac{3H^2}{2\varphi\dot\varphi^2}\,d\varphi \\
  &= \int\frac{3}{2\varphi}\left(\frac{H}{\dot\varphi}\right)^2 d\varphi.
\end{align}
From the Friedmann relation, $H/\dot\varphi = \sqrt{\frac{4\pi G\hbar^2}{3c^2M}}$ is constant. Hence
\begin{equation}
N = \frac{3}{2}\left(\frac{4\pi G\hbar^2}{3c^2M}\right)\int\frac{d\varphi}{\varphi}
    = \frac{2\pi G\hbar^2}{c^2M}\,\ln\frac{\varphi_f}{\varphi_i}.
\end{equation}
Thus $N$ is logarithmic in $\varphi$, exactly as in the Starobinsky model. A detailed perturbation analysis (see Chapter~5 of the QGD monograph) then yields the standard Starobinsky values
\begin{equation}
n_s = 1 - \frac{2}{N},\qquad r = \frac{12}{N^2},
\end{equation}
which for $N=60$ give $n_s\approx0.967$ and $r\approx0.003$, matching the Planck~2018 results.

\subsection{Conclusion}
The $\sigma$-kinetic term together with the fourth‑order stiffness term generates an effective $R^2$ term in the cosmological action. The massive mode of the $\sigma$-field plays the role of the inflaton, and the slow‑roll dynamics lead precisely to the Starobinsky inflationary predictions, which are in excellent agreement with current observations. QGD therefore naturally incorporates inflation without the need for an additional scalar field.
